\documentclass{article}
\usepackage{amsmath, amssymb, amsthm}
\usepackage{hyperref}
\usepackage{geometry}
\geometry{margin=1in}

\title{Erd\H{o}s Problem \#218}
\date{Last updated: 2025-08-31}
\author{Source: erdosproblems.com}

\begin{document}
\maketitle

\noindent\textbf{Status:} OPEN \quad \textbf{No prize}

\noindent\textbf{Tags:} number theory, primes

\medskip
\noindent\textbf{Problem Statement:}

Let $d_n=p_{n+1}-p_n$. The set of $n$ such that $d_{n+1}\geq d_n$ has density $1/2$, and similarly for $d_{n+1}\leq d_n$. Furthermore, there are infinitely many $n$ such that $d_{n+1}=d_n$.

\bigskip
\noindent\textbf{Remarks:}

In \cite{Er85c} Erd\H{o}s also conjectures that $d_n=d_{n+1}=\cdots=d_{n+k}$ is solvable for every $k$ (which is equivalent to $k$ consecutive primes in arithmetic progression, see [141]).

Banks \cite{Ba23} has given a heuristic argument, assuming a quantitative form of the prime tuples conjecture, that the first statement is true, and in fact for any $c\geq 0$ and $\epsilon>0$ the number of $n$ such that $p_n\leq x$ and $d_{n+1}\geq cd_n$ is\[\frac{\pi(x)}{c+1}+O\left(\frac{x}{(\log x)^{3/2+\epsilon}}\right).\]

\bigskip
\noindent\textbf{References:}

[Ba23] Banks, William D., On ratios of consecutive prime gaps. Integers (2023), Paper No. A50, 13.
 
 [Er85c] Erd\H{o}s, P., On some of my problems in number theory I would most like to see solved. Number theory (Ootacamund, 1984) (1985), 74-84.

\bigskip
\noindent\small{Source: \url{https://www.erdosproblems.com/218}}
\end{document}
